\begin{helper}
Use the paired candidate events of
\noderef{SamplingSplitOutsideBlockerPairedCandidateEvent}. Assume the
blocker list and private incidence map are injective, every listed blocker
lies outside \(X\cup\{r\}\), and every private vertex is distinct from
every \(\phi(y)\) with \(y\in X\cup\{r\}\). Fix \(k<n\).
Suppose paired states \(\zeta,\xi\) have identical activation indicators
and priorities at every source coordinate except possibly \(z_k\), and
at every target coordinate except possibly the vertices \(\beta(x,z_k)\)
for valid incidences with \(x\in X\). Then for every \(x\in X\),
\[
 \zeta\in F_{\{i:i<k\},\{i:k<i\}}(x)
 \quad\Longleftrightarrow\quad
 \xi\in F_{\{i:i<k\},\{i:k<i\}}(x).
\]
This asserts coordinate invariance, not probabilistic independence.
\end{helper}
\begin{proof}
Expand \noderef{SamplingSplitOutsideBlockerPairedCandidateEvent}.
The source candidate and embedded-neighbor coordinates lie in
\(X\cup\{r\}\), so none equals \(z_k\). All their activation and
priority tests agree between the two states. The corresponding target
coordinates \(\phi(y)\) avoid every current private vertex by freshness,
so their activation and priority tests also agree.

For a retained source test with \(k<i\), injectivity of the blocker list
gives \(z_i\ne z_k\). Its activation, blocker priority, and candidate
priority therefore agree. For a private test with \(i<k\), suppose its
coordinate \(\beta(x,z_i)\) equalled some current coordinate
\(\beta(a,z_k)\). Injectivity of \(\beta\) would identify the two
incidences and hence give \(z_i=z_k\); injectivity of the list would then
give \(i=k\), a contradiction. Thus this private coordinate is also
unchanged. Every atomic clause in the candidate event has the same truth
value in the two states. Conjunction and the displayed universal
quantifiers preserve this equivalence, proving the assertion for each
fixed candidate.
\end{proof}
